\section*{अध्याय \ref{ch:b1:field-lab-methods-1} --- क्षेत्र और प्रयोगशाला विधियाँ I}

\begin{solution}{exo:b1:field-lab-methods-1:1}
1:10\,000 पर \qty{4.9}{m} $4.9/10\,000 = \qty{0.49}{mm}$ हो जाता है; 1:25\,000 पर $4.9/25\,000 = \qty{0.20}{mm}$,
लगभग एक पेंसिल की रेखा की चौड़ाई।
\end{solution}

\begin{solution}{exo:b1:field-lab-methods-1:2}
रेक्याविक पर $D = \ang{-11.28}$: $52 - 11.28 = \ang{040.7}$। वेलिंगटन पर $D = \ang{+23.23}$: $52 + 23.23 = \ang{075.2}$।
वही दिक्सूचक पाठ \ang{34.5} दूर की दिशाएँ बताता है।
\end{solution}

\begin{solution}{exo:b1:field-lab-methods-1:3}
$20\to015$। $t = \ang{15}$ और $p = \ang{20}$ के साथ नीचे की ओर इकाई सदिश है
\[
  (\cos\ang{20}\sin\ang{15},\ \cos\ang{20}\cos\ang{15},\ -\sin\ang{20}) =
  (0.2432,\ 0.9077,\ -0.3420).
\]
\end{solution}

\begin{solution}{exo:b1:field-lab-methods-1:4}
अंकगणितीय माध्य $(350 + 358 + 6)/3 = \ang{238}$ है, जो तीनों पाठों से लगभग विपरीत है। क्षैतिज इकाई सदिश
$(\sin a, \cos a)$ जुड़कर $(-0.1736 - 0.0349 + 0.1045,\ 0.9848 + 0.9994 + 0.9945) = (-0.1040, 2.9787)$ देते हैं,
जिनका \omterm{def:b1:field-lab-methods-1:trend-plunge}{दिगंश} $\operatorname{atan2}(-0.1040, 2.9787) = \ang{-2.0}$ है, यानी \ang{358.0}: सही माध्य, पाठों के बीच में।
\end{solution}

\begin{solution}{exo:b1:field-lab-methods-1:5}
माध्य $175/5 = \ang{35.0}$। विचलन $-1, 2, 0, -2, 1$; वर्गों का योग 10;
$s = \sqrt{10/4} = \ang{1.58}$। मानक त्रुटि $1.58/\sqrt5 = \ang{0.71}$।
\end{solution}

\begin{solution}{exo:b1:field-lab-methods-1:6}
अभिलंब $(\sin\ang{30}, 0, \cos\ang{30}) = (0.5, 0, 0.866)$ और
$(-0.5, 0, 0.866)$ हैं; $\vec S = (0, 0, 1.732)$ ऊर्ध्वाधर है: माध्य एक क्षैतिज समतल है, और $\bar R = 1.732/2 = 0.866$, हर
पाठ माध्य से \ang{30} दूर। दोनों पाठ एक सतह को त्रुटियों के साथ प्रतिचयित नहीं करते: वे एक \omterm{def:b1:geological-maps:fold}{वलन} की दो भुजाओं के हैं,
दो भिन्न समष्टियाँ, और उनका औसत न एक का वर्णन करता है न दूसरे का। हर भुजा अपने आप औसत की जाती है; दोनों मिलकर
\omterm{def:b1:geological-maps:fold}{वलन-अक्ष} तय करते हैं, दोनों समतलों में पड़ी रेखा, यहाँ क्षैतिज और उत्तर-दक्षिण।
\end{solution}

\begin{solution}{exo:b1:field-lab-methods-1:7}
$\tan p = \tan\ang{12}\cos(\ang{80} - \ang{120}) = 0.2126\times0.7660 =
0.1628$: $p = \ang{9.2}$, खरोंचें \ang{120} की ओर नमित हैं। \omterm{def:b1:field-lab-methods-1:trend-plunge}{दिगंश} \ang{170} के लिए, $\cos(\ang{-90}) = 0$: रेखा क्षैतिज है,
\omterm{def:b1:geological-maps:attitude}{स्ट्राइक} के अनुदिश।
\end{solution}

\begin{solution}{exo:b1:field-lab-methods-1:8}
$0.47/0.20 = 2.35$ वर्ष 1 जनवरी 2026 के बाद, यानी वसंत 2028 में। 1 जनवरी 2030 को:
$-0.47 + 4\times0.20 = \ang{+0.33}$, भौगोलिक उत्तर के पूर्व में।
\end{solution}

\begin{solution}{exo:b1:field-lab-methods-1:9}
तीनों अभिलंब \omterm{def:b1:field-lab-methods-1:trend-plunge}{दिगंश} \ang{080} के ऊर्ध्वाधर तल में हैं, ऊर्ध्वाधर से \ang{10}, \ang{12} और \ang{14} पर; बाहरी दो बीच वाले के
सापेक्ष सममित हैं, अतः उनका योग उसी के अनुदिश है: माध्य $080/12$ है, $\theta = \ang{2}, 0, \ang{2}$ के साथ। तब
$\bar R = (1 +
2\cos\ang{2})/3 = (1 + 1.998782)/3 = 0.999594$, और $\sqrt{2(1 - \bar R)} =
\sqrt{0.000812} = \qty{0.02850}{rad} = \ang{1.633}$; ठीक-ठीक,
$\theta_{\text{वर्ग-माध्य-मूल}} = \sqrt{(4 + 0 + 4)/3} = \ang{1.633}$।
\end{solution}

\begin{solution}{exo:b1:field-lab-methods-1:10}
काट त्रिज्या $\sqrt{r^2 - h^2}$ का एक वृत्त है। उसका माध्य
$\frac1r\int_0^r\sqrt{r^2 - h^2}\,\mathrm dh = \frac1r\cdot\frac{\pi r^2}{4}
= \frac{\pi r}{4}$ है, क्योंकि समाकल एक चौथाई चकती का क्षेत्रफल है। \qty{0.50}{mm} आर-पार के कण औसतन
$0.50\times\pi/4 = \qty{0.39}{mm}$ आर-पार दिखते हैं: \omterm{def:b1:field-lab-methods-1:thin-section}{पतली काट} कण-आकारों को कम आँकती है।
\end{solution}

\begin{solution}{exo:b1:field-lab-methods-1:11}
ऊर्ध्वाधर के परितः $D$ का घूर्णन, यानी
$(x, y, z) \mapsto (x\cos D + y\sin D,\ -x\sin D + y\cos D,\ z)$, भेजता है
$(\sin\delta\sin\alpha, \sin\delta\cos\alpha, \cos\delta)$ को
$(\sin\delta\sin(\alpha + D), \sin\delta\cos(\alpha + D), \cos\delta)$ में: यह हर \omterm{def:b1:geological-maps:attitude}{नति-दिशा} में $D$ जोड़ता है। रैखिक होने
के कारण, यह परिणामी $\vec S$ को घुमाए गए अभिलंबों के परिणामी में भेजता है; यह लंबाइयों और ऊर्ध्वाधर घटक को बनाए
रखता है, अतः माध्य \omterm{def:b1:geological-maps:attitude}{नति} $\arccos(S_U/\|\vec S\|)$ अपरिवर्तित रहती है, और यह $\vec S$ के क्षैतिज भाग को $D$ घुमाता है,
इसलिए माध्य \omterm{def:b1:geological-maps:attitude}{नति-दिशा} $D$ बढ़ती है। हर पाठ को विचलन के लिए सुधारने पर, या केवल उनके माध्य को, वही परिणाम मिलता है।
\end{solution}

\begin{solution}{exo:b1:field-lab-methods-1:12}
सबसे अधिक \omterm{def:b1:geological-maps:attitude}{नति} की रेखा के अनुदिश प्रकीर्णन \omterm{def:b1:geological-maps:attitude}{नति} का प्रकीर्णन है: मानक त्रुटि $2/\sqrt N \le 0.5$ है, $N \ge 16$ के लिए।
उसके आर-पार, अभिलंब का \ang{2} घुमाव \omterm{def:b1:geological-maps:attitude}{नति-दिशा} को $2/\sin\ang{12} = \ang{9.62}$ बदलता है; माध्य \omterm{def:b1:geological-maps:attitude}{नति-दिशा} की मानक
त्रुटि $9.62/\sqrt N \le 1$ है, $N \ge 92.5$ के लिए, यानी 93 पाठ। मंद \omterm{def:b1:geological-maps:attitude}{नति} जल्दी ज्ञात होती है; उसकी दिशा नहीं।
\end{solution}

\begin{solution}{pb:b1:field-lab-methods-1:1}
\textbf{1.} नति-दिशाएँ: $967/12 = \ang{80.6}$; नतियाँ: $145/12 =
\ang{12.1}$।

\textbf{2.} \ang{355} और \ang{005} का औसत \ang{180} है, विपरीत दिशा। यहाँ सारी नति-दिशाएँ \ang{076} और \ang{085} के बीच हैं,
\ang{000} पर की छलांग से दूर, इसलिए अंकगणितीय माध्य काम करता है।

\textbf{3.} जाँच:
\[ \|\vec n\|^2 = \sin^2\delta(\sin^2\alpha + \cos^2\alpha) + \cos^2\delta = 1 \]
नीचे की ओर नति-रेखा की दिशा का सदिश, यानी
$(\cos\delta\sin\alpha, \cos\delta\cos\alpha, -\sin\delta)$ है, और $\vec n$ से उसका अदिश गुणनफल
$\sin\delta\cos\delta - \cos\delta\sin\delta = 0$ है; स्ट्राइक-रेखा क्षैतिज और \omterm{def:b1:geological-maps:attitude}{नति-दिशा} के लंबवत है,
$(\cos\alpha, -\sin\alpha, 0)$, और $\vec n$ से उसका अदिश गुणनफल
$\sin\delta\sin\alpha\cos\alpha - \sin\delta\cos\alpha\sin\alpha = 0$ है।

\textbf{4.} पाठ 11, पाठ 2 के समान है: $(0.2222, 0.0352, 0.9744)$। पाठ 12:
\[
  (\sin\ang{11}\sin\ang{82},\ \sin\ang{11}\cos\ang{82},\ \cos\ang{11}) =
  (0.1890,\ 0.0266,\ 0.9816).
\]

\textbf{5.} $\vec S = (2.4747,\ 0.4082,\ 11.7294)$।

\textbf{6.} $\|\vec S\| = \sqrt{6.1241 + 0.1666 + 137.5788} = 11.99457$, और
$\bar R = 11.99457/12 = 0.999547$।

\textbf{7.} $\bar\alpha = \operatorname{atan2}(2.4747, 0.4082) = \ang{80.6}$
और $\bar\delta = \arccos(11.7294/11.99457) = \arccos 0.977893 = \ang{12.1}$:
सदिश माध्य $080.6/12.1$ है, यहाँ के अंकगणितीय माध्यों के समान।

\textbf{8.} नति-दिशाएँ: $s = \ang{2.94}$; नतियाँ: $s = \ang{1.68}$।

\textbf{9.} $\theta_{\text{वर्ग-माध्य-मूल}} \approx \sqrt{2\times0.000453} =
\qty{0.0301}{rad} = \ang{1.72}$।

\textbf{10.} $\sqrt{1.68^2 + (\sin\ang{12.1}\times2.94)^2} = \sqrt{2.82 +
0.38} = \ang{1.79}$, \ang{1.72} के निकट: पाठ त्रुटि-प्रतिरूप के अनुसार व्यवहार करते हैं, एक अंश तक पूर्णांकन के प्रभावों
और एक छोटे नमूने के भीतर।

\textbf{11.} \omterm{def:b1:geological-maps:attitude}{नति}: वह 2.82 (वर्ग-अंशों में) का 3.20 में योगदान करती है। \omterm{def:b1:geological-maps:attitude}{नति-दिशा} केवल $\sin\delta = 0.21$ के
माध्यम से गिनी जाती है, क्योंकि मंद शय्या पर \omterm{def:b1:geological-maps:attitude}{नति-दिशा} का बदलाव समतल को बहुत कम घुमाता है।

\textbf{12.} माध्य \omterm{def:b1:geological-maps:attitude}{नति}: $1.68/\sqrt{12} = \ang{0.48}$; माध्य \omterm{def:b1:geological-maps:attitude}{नति-दिशा}:
$2.94/\sqrt{12} = \ang{0.85}$।

\textbf{13.} नहीं: \ang{12} \omterm{def:b1:geological-maps:attitude}{नति} वाली शय्या पर \omterm{def:b1:geological-maps:attitude}{नति-दिशा} के \ang{2.9} का अर्थ केवल
$2.9\sin\ang{12.1} = \ang{0.6}$ समतल का घुमाव है, जो नतियों के \ang{1.7} से कम है। मंद शय्या-सतह का कोई भी हल्का
उभार उसकी \omterm{def:b1:geological-maps:attitude}{नति-दिशा} को दूर तक बिखेर देता है।

\textbf{14.} $D = \ang{-2.5}$: $83.1 - 2.5 = \ang{080.6}$; उसका माध्य $080.6/12.1$ है, पहले भूवैज्ञानिक के समान।

\textbf{15.} \omterm{def:b1:geological-maps:attitude}{नति} को ग्रुत्वाकर्षण द्वारा दिए गए ऊर्ध्वाधर के सापेक्ष क्लिनोमीटर से पढ़ा जाता है; सुई की कोई भूमिका नहीं।

\textbf{16.} हाँ: ऊर्ध्वाधर के परितः घूर्णन हर पाठ और माध्य को एक ही कोण से हिलाता है (\cref{exo:b1:field-lab-methods-1:11});
उत्तर वही है।

\textbf{17.} हर \omterm{def:b1:geological-maps:attitude}{नति-दिशा} पर $10\times0.2 = \ang{2}$ की त्रुटि, \ang{0.85} की मानक त्रुटि से दुगने से अधिक: एक व्यवस्थित
त्रुटि जिसे औसत नहीं पकड़ सकता।

\textbf{18.} $80.6 - 90 = \ang{350.6}$; \omterm{def:b1:geological-maps:attitude}{अभिविन्यास} \omterm{def:b1:geological-maps:attitude}{स्ट्राइक} \ang{351}, \omterm{def:b1:geological-maps:attitude}{नति} \ang{12} पूर्व की ओर।

\textbf{19.} नति-दिशाएँ $080.6$ और $080.6$: इस परिशुद्धता पर कोई अंतर नहीं; नतियाँ $12.1$ और $12.2$: \ang{0.1}।

\textbf{20.} अभिलंब $(\sin\ang{12.2}\sin\ang{80.6},
\sin\ang{12.2}\cos\ang{80.6}, \cos\ang{12.2})$ और $\vec S/\|\vec S\|$ \ang{0.13} का कोण बनाते हैं।

\textbf{21.} $1.72/\sqrt{12} = \ang{0.50}$; अंतर, \ang{0.13}, उसका एक चौथाई है: सार्थक नहीं। कम्पास और बोरहोल सहमत हैं।

\textbf{22.} चार गुना अधिक: 48 पाठ।

\textbf{23.} उदाहरण: एक गलत या पुरानी विचलन-सेटिंग; पास के लोहे (वाहन, औज़ार, चुंबकीय चट्टानें) से विक्षेपित दिक्सूचक; एक
बिगड़ा क्लिनोमीटर; केवल सबसे अच्छी तरह खुली सतहों पर लिए गए पाठ, जो भुजा का प्रतिनिधित्व न करें। औसत यादृच्छिक
त्रुटियाँ घटाता है, व्यवस्थित त्रुटियाँ नहीं।

\textbf{24.} पूर्वी भुजा का माध्य \omterm{def:b1:geological-maps:attitude}{अभिविन्यास} \textbf{081/12} है: \omterm{def:b1:geological-maps:attitude}{नति-दिशा} \ang{081}, \omterm{def:b1:geological-maps:attitude}{नति} \ang{12}।
\end{solution}
